% ======================================================================
%  TIER 2 KARTI: Veri & Gizlilik
% ======================================================================

\card{Veri \& Gizlilik}{verin nereye gidiyor, nasıl korunuyor}
\cl{Problem} Bir AI özelliği kurarken en kritik sorular çoğu zaman teknik değil,
\emph{veriyle} ilgilidir: ``Kullanıcı verimiz modelin eğitiminde kullanılıyor mu?
Nerede işleniyor/saklanıyor? Sisteme beslediğimiz veri yeterince temiz mi?''
Yanlış cevap; yasal risk + kötü sonuç demektir.
\cl{Tanım} İki ayrı kaygı: \textbf{gizlilik/uyum} --- verinin kullanımı, saklanması
ve coğrafyası (\term{data residency}, GDPR); ve \textbf{veri kalitesi} --- sisteme
beslenen verinin doğruluğu/tazeliği (\en{garbage in, garbage out}).
\cl{Model} Model bir mutfak, verin malzeme. Kötü malzemeyle iyi yemek çıkmaz ---
ve malzemenin nereden geldiği, nereye gittiği denetlenebilir olmalıdır.
\cl{PM soruları} Sağlayıcı verimizi eğitimde kullanıyor mu (\en{zero-retention}/opt-out var mı)?
Veri hangi bölgede işleniyor? PII maskeleniyor mu? RAG kaynağı güncel ve doğru mu?
\cl{Data flywheel} Production kullanımı (gerçek sorular, düzeltmeler) → yarının
eval seti ve iyileştirme yakıtı. En değerli veri kullanımdan doğar.
\cl{Tuzak} Veriyi sonradan düşünmek; sağlayıcı sözleşmesini okumamak; kirli veriyle
RAG kurup ``model kötü'' demek.
\cl{İlişkili} \term{RAG} · \term{Fine-tuning} · \term{Guardrails} (PII) · \term{Eval} (flywheel) \,{\scriptsize\color{softgray}(İlke \#18, \#22)}
